\section{估值门槛与本案的正式降级}
完整估值至少需要同一日期的价格、股本、市值、收入/利润/EPS 分母、业务模型匹配的方法、熊基牛情景、市场隐含预期及可用的卖方比较。本案缺少其中多项关键输入，且供应链审计未选出 eligible/core 公司。因此不产生假精确的 PE、目标价或上涨空间，正式降级为证据型研究报告。

\begin{exhibitbox}[升级与否定条件]
\small
\begin{tabularx}{\exhibitboxwidth}{L{2.1cm}X X}
\toprule 方向 & 升级为估值候选的条件 & 主要否定/降级信号\\ \midrule
半导体设备 & 客户/晶圆厂认证、订单或验收、产品 ASP、毛利和现金流同时可核。 & 需求仅停留在海外参照或行业热度；订单/利用率未兑现。\\
创新药 BD & 逐资产会计、里程碑概率、成本承担、海外商业化和 EPS 桥可核。 & 把潜在里程碑、合作方名称或一次确认视为常态利润。\\
电力设备 & 中标份额、订单、交付验收、价格成本、应收和现金转换可核。 & 政策/试点/装机替代订单与回款；竞争降价侵蚀毛利。\\
\bottomrule
\end{tabularx}
\end{exhibitbox}

\section{风险矩阵}
除基本面风险外，本文特别强调证据风险：文章三项资金数字不可复算、样本当前价格和财务口径未完整冻结、卖方目标覆盖不足。证据风险不是脚注；它直接决定不能给出核心仓位和估值信用。
